%% Chapter/chapter3.tex
%% บทที่ 3 - วิธีการดำเนินการ

\chapter{วิธีการดำเนินการ}
% (English: Methodology)

\noindent\hspace{\parindent}%
[กล่าวนำถึงวิธีการดำเนินการวิจัย อธิบายภาพรวมของระบบหรือแบบจำลองที่ใช้ในการวิจัย]

%%------------------------------------------------------------------
\section{แบบจำลองหรือระบบ}

\noindent\hspace{\parindent}%
[อธิบายแบบจำลองหรือระบบที่ใช้ในการวิจัย รวมถึงค่าพารามิเตอร์ต่าง ๆ]

ตัวอย่างตาราง:
\begin{table}[htbp]
  \centering
  \caption{ชื่อตาราง: ค่าพารามิเตอร์ของระบบ}
  \label{tab:parameters}
  \begin{tabular}{c|l|c|c}
    \toprule
    สัญลักษณ์ & คำอธิบาย & ค่า & หน่วย \\
    \midrule
    $T_G$   & ค่าคงที่เวลาของ Governor     & 0.08  & วินาที \\
    $T_T$   & ค่าคงที่เวลาของ Turbine      & 0.30  & วินาที \\
    $T_R$   & ค่าคงที่เวลาของ Reheat       & 10.0  & วินาที \\
    $K_R$   & อัตราขยายของ Reheat          & 0.5   & -- \\
    $K_P$   & อัตราขยายของระบบกำลัง        & 120   & Hz/p.u. \\
    $T_P$   & ค่าคงที่เวลาของระบบกำลัง     & 20.0  & วินาที \\
    \bottomrule
  \end{tabular}
\end{table}


%%------------------------------------------------------------------
\section{ขอบเขตการค้นหาพารามิเตอร์}

\noindent\hspace{\parindent}%
[อธิบายขอบเขตของค่าพารามิเตอร์ที่ใช้ในการหาค่าที่เหมาะสม]

%ในการทดลองนี้ได้กำหนดขอบเขตของค่าพารามิเตอร์ไว้ดังสมการ~\eqref{eq:bounds}

%\begin{equation}
%  \mathbf{x}_{min} \leq \mathbf{x} \leq \mathbf{x}_{max}
%  \label{eq:bounds}
%\end{equation}

%%------------------------------------------------------------------
\section{ฟังก์ชันวัตถุประสงค์}

\noindent\hspace{\parindent}%
อธิบายฟังก์ชันวัตถุประสงค์ที่ใช้ในการประเมินผลลัพธ์

เกณฑ์ดัชนีสมรรถนะที่ใช้ในการหาค่าที่เหมาะสม:

%\begin{equation}
%  \text{ITAE} = \int_{0}^{t_{sim}} t \left| \Delta f_i + \Delta P_{tie,i} \right| dt
%  \label{eq:itae}
%\end{equation}

%\begin{equation}
%  \text{ITSE} = \int_{0}^{t_{sim}} t \left[ (\Delta f_i)^2 + (\Delta P_{tie,i})^2 \right] dt
%  \label{eq:itse}
%\end{equation}

%\begin{equation}
%  \text{IAE} = \int_{0}^{t_{sim}} \left| \Delta f_i + \Delta P_{tie,i} \right| dt
%  \label{eq:iae}
%\end{equation}

%\begin{equation}
%  \text{ISE} = \int_{0}^{t_{sim}} \left[ (\Delta f_i)^2 + (\Delta P_{tie,i})^2 \right] dt
%  \label{eq:ise}
%\end{equation}

%%------------------------------------------------------------------
\section{สถานการณ์จำลอง}

\noindent\hspace{\parindent}%
[อธิบายสถานการณ์การจำลองที่ใช้ในการทดสอบ]

การทดลองนี้แบ่งออกเป็น 2 กรณีศึกษา ดังนี้
\begin{enumerate}
  \item กรณีศึกษาที่ 1 อธิบายโหลดโปรไฟล์แบบที่ 1
  \item กรณีศึกษาที่ 2 อธิบายโหลดโปรไฟล์แบบที่ 2
\end{enumerate}

%%------------------------------------------------------------------
\section{การวิเคราะห์ผลลัพธ์}

\noindent\hspace{\parindent}%
[อธิบายวิธีการวิเคราะห์ผลลัพธ์และตัวชี้วัดที่ใช้ในการประเมิน]

\begin{itemize}
  \item \textbf{Maximum Overshoot}: ค่าความเบี่ยงเบนสูงสุดในทิศทางบวก
  \item \textbf{Minimum Undershoot}: ค่าความเบี่ยงเบนต่ำสุดในทิศทางลบ
  \item \textbf{Settling Time}: ระยะเวลาที่ระบบเข้าสู่สภาวะคงตัว
  \item \textbf{Steady-State Error}: ความผิดพลาดในสภาวะคงตัว
\end{itemize}
